% !PRINTED 413
\weberpage{413}
Nehmen wir hierin $\alpha=\alpha_1$ und $\alpha=\alpha_2$ an, und setzen zur
Abkürzung $\varphi_{\alpha_1}=\varphi_1$, $\varphi_{\alpha_2}=\varphi_2$, so
können wir aus den beiden Gleichungen
\[
\varphi_1 = a_1\varphi + b_1\psi, \quad \varphi_2 = a_2\varphi + b_2\psi
\]
$\psi$ eliminiren und erhalten eine Gleichung von der Form:
\weq{7}{h\varphi + h_1\varphi_1 + h_2\varphi_2 = 0,}
worin $h$, $h_1$, $h_2$ Constanten sind, unter denen wenigstens zwei von
Null verschieden sind. Auf die identische Gleichung (7) können wir nun die
Permutationen $\alpha_1$, $\alpha_2$ anwenden, und erhalten, da
$\alpha_1\alpha_2 = \alpha_2\alpha_1 = \alpha_3$, $\alpha_1\alpha_1 = 1$,
$\alpha_2\alpha_2 = 1$ ist,
\[
\begin{array}{l|l}
h\varphi + h_1\varphi_1 + h_2\varphi_2 = 0 & h,\\
h\varphi_1 + h_1\varphi + h_2\varphi_3 = 0 & h_1,\\
h\varphi_2 + h_1\varphi_3 + h_2\varphi = 0 & -h_2,
\end{array}
\]
woraus durch Multiplication mit den daneben stehenden Factoren und
Addition:
\weq{8}{(h^2 + h_1^2 - h_2^2)\varphi + 2hh_1\varphi_1 = 0.}
Wenn also $h$ und $h_1$ von Null verschieden sind, so unterscheiden sich
$\varphi$ und $\varphi_1$ nur durch einen constanten Factor (der $=\pm 1$
sein muss, da $\alpha_1$ vom 2\textsuperscript{ten} Grade ist). Ist aber
$h$ oder $h_1=0$, so folgt aus (7) $\varphi_1=\pm\varphi_2$ oder
$\varphi=\pm\varphi_2$, und die Anwendung von $\alpha_1$ auf die erste
Gleichung giebt $\varphi=\pm\varphi_3$. Es findet also jedenfalls \gesperrt{eine der
Relationen} statt:
\weq{9}{\varphi = \pm\varphi_1, \quad \varphi = \pm\varphi_2, \quad \varphi = \pm\varphi_3.}
Da man nun ebensowohl $\gamma,\alpha_1$ als $\gamma,\alpha_2$, als auch
$\gamma,\alpha_3$ als erzeugende Elemente der Gruppe $A$ betrachten kann,
so ergiebt dies Resultat, zusammengenommen mit (4), den Satz:

1. \gesperrt{Die Function $\varphi$ hat die Eigenschaft, sich durch alle
Permutationen der alternirenden Gruppe nur um einen constanten Factor zu
ändern.}

Ganz derselbe Schluss ist aber auch auf $\psi$ anwendbar, und also auch
auf den Quotienten $\vartheta$ beider Functionen. Es ist sonach für jede
Permutation $\pi$ der alternirenden Gruppe
\weq{10}{\vartheta_\pi = \varepsilon\vartheta.}
Hierin ist, wenn $\pi$ zu den cyklischen Permutationen $\gamma$ gehört,
$\varepsilon$ eine dritte Einheitswurzel. Daraus ist aber ferner zu
schliessen, weil alle Permutationen $\pi$ aus $\gamma$, $\gamma'$
zusammengesetzt
%CONT>>
